\section{The Module Definition}\label{sec:ModuleDefinition}
Modules\index{module} have been introduced to Java with version~9, as a result of the Jigsaw project\index{Project Jigsaw|see {Java Platform Module System}}\index{Java Platform Module System} \autocite{OPENJDK:ProjectJigsaw}, today mostly referenced as JPMS.

A module will be defined in the file \verb#module-info.java#\index{module-info.java} that is located in the root of the source tree for the project. Details can be found in \autocite{ORACLE_DOC_LANGUAGE_SPECIFICATION:ModuleDeclarations}. According to \autocite{ORACLE_DOC_LANGUAGE_SPECIFICATION:CompilationUnit} is the \verb#module-info.java#\index{module-info.java} file a special kind of a compilation unit\index{compilation unit}: a \bdq{Modular Compilation Unit}. Basically, that file will look like this:
\begin{lstlisting}[numbers=left, caption={module-info.java},morekeywords={exports,module,opens,provides,requires,to,transitive,uses,var,with,yield}]
/*
 * ==================================================================
 * Copyright © <Year> <Copyright Notice>
 * ==================================================================
 *
 * <License Notice>
 */
 
/**
 *  <module description>
 */ 
module <modulename> 
{
    requires <name_of_required_module_1>;
    requires <name_of_required_module_2>;

    requires transitive <name_of_required_module_3>;

    exports <name_of_exported_package_1>;
    exports <name_of_exported_package_2> to <name_of_target_module_1>;
(|\newpage|)
    opens <name_of_opened_package_1>;
    opens <name_of_opened_package_1> to <name_of_target_module_1>, <name_of_target_module_2>;

    uses <name_of_provider_interface>;
    provides <name_of_provider_interface> with <name_of_provider_implementation>;
}
//  module <modulename>

/*
 *  End of File
 */
\end{lstlisting}
How to define the name of the module (\lstinline|<modulename>|) is covered in chapter \tqfullvref{sec:Modules}, modularisation\index{modularisation} in general (and the contents of a \verb#module-info.java#\index{module-info.java} in detail) is discussed in \tqfullvref{sec:Modularisation}.

The \verb#<module description># will be discussed in chapter \tqfullvref{sec:ModuleComment}.

\subsubsection{Principles}
\begin{enumerate}[label=P\arabic*.]
	\item{A \verb#module-info.java#\index{module-info.java} has a beginning and a closing comment as a regular Java source file. See \tqfullref{sec:BeginningComment} and \tqfullref{sec:ClosingComment}.}
	\item{It is possible that a \verb#module-info.java#\index{module-info.java} will import something.

The \lstinline|import| statements follow the beginning comment, separated with an empty line.}
	\item{The module description provides brief information about the purpose of the module.}
	\item{Following the module definition block, there has to be a comment that repeats the module name.}
	\item{Each keyword \lstinline|requires|, \lstinline|exports|, \lstinline|opens|, \lstinline|uses| and \lstinline|provides| has its own line.
	
	It will be indented one level.}
	\item{The keywords will be grouped, and the groups are separated by blank lines.
	
	If a \lstinline|uses| and a \lstinline|provides| clause refer to the same provider, they should be placed together (without a blank line).}
\end{enumerate}

%==============================================================================
% End of file!!
%==============================================================================
